\documentclass[11pt,a4paper]{article}
\usepackage{amsmath,amssymb}
\usepackage{geometry}
\usepackage{hyperref}
\usepackage{booktabs}
\usepackage{listings}
\usepackage{xcolor}
\geometry{margin=2.5cm}

\title{FRED 2.0 Platform\\[6pt]
       \large Theory and Developer Manual}
\author{FRED Development Team}
\date{\today}

\lstset{
  basicstyle=\small\ttfamily,
  keywordstyle=\color{blue},
  commentstyle=\color{gray},
  breaklines=true,
  frame=single
}

\begin{document}
\maketitle
\tableofcontents
\newpage

%===========================================================================
\section{Introduction}
%===========================================================================

FRED 2.0 is a modular, open-source platform for developing fuel performance
codes for nuclear fuel rods.  The platform is written in C\texttt{++} and
provides a Python interface via \texttt{pybind11} so that all problem
set-up---geometry, material selection, boundary conditions---is performed
through Python scripts.  Time integration is performed by the SUNDIALS IDA
solver (Implicit Differential-Algebraic equations).

The platform targets cylindrical fuel rod geometries with radial and axial
discretisation.  Three applications are built on top of the platform:

\begin{description}
  \item[\textbf{FRED-ROD}] Pure thermo-mechanical solver.  No irradiation
        model.  Suitable for out-of-pile thermal tests, gap-closure studies,
        and code verification.
  \item[\textbf{FRED-OX}]  Mixed oxide (MOX) fuel in a sodium-cooled
        environment, including fission gas release, fuel swelling, and
        irradiation creep.
  \item[\textbf{FRED-M-Na}]  Metallic fuel (U-Pu-Zr) in sodium.
        Includes irradiation-specific phenomena for metallic fuel.
\end{description}

%===========================================================================
\section{Code Architecture}
%===========================================================================

\subsection{Directory Structure}

\begin{lstlisting}
fred_platform/
  src/
    materials/              # Material property interfaces
      FuelPelletMaterial.hpp     # Abstract base: fuel
      CladdingMaterial.hpp       # Abstract base: cladding
      CoolantMaterial.hpp        # Abstract base: coolant
      uo2/  UO2.hpp/.cpp         # UO2 concrete implementation
      aim1/ AIM1.hpp/.cpp        # AIM1 cladding
      na/   Na.hpp/.cpp          # Liquid sodium
    platform/               # Shared solver components
      Constants.hpp              # Physical constants
      FuelRodGeometry.hpp        # Geometry data structure
      FuelRodState.hpp           # Runtime state (temperatures, stresses)
      GapModel.hpp/.cpp          # Fuel-clad gap conductance
      HeatConduction.hpp/.cpp    # Thermal residuals
      StressStrain.hpp/.cpp      # Mechanical residuals
    apps/
      fred_rod/
        FredRodResiduals.hpp/.cpp  # Full DAE residual (thermal + mechanical)
        FredRodSolver.hpp/.cpp     # SUNDIALS IDA driver
        bindings.cpp               # pybind11 Python module
        main.cpp                   # Stand-alone C++ driver
  doc/
    fred_platform.tex       # This document
  Makefile
\end{lstlisting}

\subsection{Material Hierarchy}

All material properties are accessed through abstract virtual interfaces:

\begin{itemize}
  \item \texttt{FuelPelletMaterial} — thermal conductivity, heat capacity,
        thermal expansion strain, Young's modulus, Poisson's ratio,
        reference/theoretical density, melting temperature.
  \item \texttt{CladdingMaterial} — same thermal and elastic properties plus
        Meyer hardness (for gap solid-contact conductance).
  \item \texttt{CoolantMaterial} — thermal conductivity, heat capacity,
        density.
\end{itemize}

Concrete implementations (\texttt{UO2}, \texttt{AIM1}, \texttt{Na}) inherit
from these interfaces.  New materials for FRED-OX and FRED-M-Na extend the
base classes to add irradiation-specific properties (swelling, fission gas,
creep) without modifying the platform layer.

%===========================================================================
\section{Mathematical Formulation}
%===========================================================================

\subsection{Cylindrical Geometry}

The fuel rod is divided into $n_z$ axial layers each of height $\Delta z$.
Each layer contains $n_f$ fuel radial nodes and $n_c$ cladding radial nodes.
The centreline of node $i$ is at radius $r_{0,i}$.  The boundary between
node $i$ and $i+1$ is at $r_{0,i+1/2}$.  Node cross-sectional areas are
\begin{equation}
  A_{0,i} = 2\pi r_{0,i} \, \Delta r,
\end{equation}
where $\Delta r = \Delta r_f$ for fuel and $\Delta r = \Delta r_c$ for the
cladding.

\subsection{Heat Conduction}

The transient energy balance for a fuel ring $i$ is
\begin{equation}
  \rho \, c_p(T_i) \, A_i \, \frac{dT_i}{dt}
    = \dot{q}(t) \, A_i + q''_{i-1/2} \cdot 2\pi r_{i-1/2}
                        - q''_{i+1/2} \cdot 2\pi r_{i+1/2},
  \label{eq:energy_ode}
\end{equation}
where $\dot{q}$ is the volumetric power density [W/m$^3$] and
\begin{equation}
  q''_{i+1/2} = k_{i+1/2}(T) \frac{T_i - T_{i+1}}{\Delta r}
\end{equation}
is the radial heat flux (positive outward).  The thermal conductivity is
evaluated at the inter-node temperature $T_{i+1/2} = (T_i + T_{i+1})/2$.

For the cladding rings no power is generated; the same form holds with
$\dot{q}=0$.

The fuel–clad gap is treated as an additional conductance:
\begin{equation}
  q''_{\mathrm{gap}} = h_{\mathrm{gap}} (T_{\mathrm{fo}} - T_{\mathrm{ci}}),
  \label{eq:gap_flux}
\end{equation}
where $T_{\mathrm{fo}}$ and $T_{\mathrm{ci}}$ are the fuel outer and cladding
inner surface temperatures, and $h_{\mathrm{gap}}$ is the total gap
conductance [W/(m$^2\cdot$K)].

The outer cladding surface temperature is prescribed as an algebraic
equation:
\begin{equation}
  T_{n_f + n_c}(t) = T_{\mathrm{coolant}}(t).
\end{equation}

\subsection{Gap Conductance Model}

The total gap conductance is the sum of three contributions (Ross \& Stoute
model):
\begin{equation}
  h_{\mathrm{gap}} = h_{\mathrm{gas}} + h_{\mathrm{rad}} + h_{\mathrm{contact}}.
  \label{eq:htot}
\end{equation}

\paragraph{Gas conductance.}
\begin{equation}
  h_{\mathrm{gas}} = \frac{k_{\mathrm{gas}}}{\delta_{\mathrm{eff}} + g_j},
  \label{eq:hgas}
\end{equation}
where $k_{\mathrm{gas}}$ is the fill-gas thermal conductivity,
$\delta_{\mathrm{eff}} = \max(\delta_{\mathrm{gap}},\,r_{\mathrm{eff}})$
is the effective gap width (bounded below by the combined surface roughness
$r_{\mathrm{eff}}=\sqrt{r_f^2 + r_c^2}$), and $g_j$ is the temperature-jump
distance (Lanning–Hann accommodation model):
\begin{equation}
  g_j = \frac{0.024688 \, k_{\mathrm{gas}} \sqrt{T_{\mathrm{avg}}}}{p_{\mathrm{gas}} \, \alpha},
\end{equation}
with accommodation coefficient $\alpha = 0.425 - 2.3\times10^{-4} \min(T_{\mathrm{avg}}, 1000)$.

Fill-gas thermal conductivity for a He/Ar/Kr/Xe mixture uses the geometric
mean mixing rule:
\begin{equation}
  k_{\mathrm{gas}} = \prod_{s} k_s^{x_s/\sum x_s},
\end{equation}
where $x_s$ are the mole fractions and $k_s(T)$ are individual gas conductivities
fit by power laws (from FRED-OX).

\paragraph{Radiation conductance.}
\begin{equation}
  h_{\mathrm{rad}} = \sigma_{\mathrm{SB}} \, F_e
    \,(T_{\mathrm{fo}}^2 + T_{\mathrm{ci}}^2)(T_{\mathrm{fo}} + T_{\mathrm{ci}}),
\end{equation}
where $F_e$ is the radiation view factor between concentric cylinders:
\begin{equation}
  F_e = \left[\frac{1}{\varepsilon_f} + \frac{r_{\mathrm{fo}}}{r_{\mathrm{ci}}}
        \left(\frac{1}{\varepsilon_c} - 1\right)\right]^{-1}.
\end{equation}
Surface emissivities: $\varepsilon_f = 0.8$, $\varepsilon_c = 0.9$
(KIT recommendation, Struwe \& Schikorr 2011).

\paragraph{Solid contact conductance (closed gap only).}
\begin{equation}
  h_{\mathrm{contact}} = \frac{33.3 \, k_m \, \Delta p_c}{H_M \sqrt{r_{\mathrm{eff}}}},
\end{equation}
where $\Delta p_c = \max(p_{\mathrm{fc}} - p_{\mathrm{gas}}, 0)$ is the net
contact pressure, $H_M$ is the Meyer hardness of the cladding, and $k_m$ is
the harmonic mean conductivity of fuel and cladding at the contact temperature.

\subsection{Thermo-Elastic Stress-Strain (FRED-ROD)}

In FRED-ROD no creep, swelling, or irradiation-induced strains are present.
The total strain is the sum of elastic and thermal contributions only:
\begin{equation}
  \varepsilon_{ij} = \varepsilon^{\mathrm{el}}_{ij} + \varepsilon^{\mathrm{th}} \delta_{ij}.
\end{equation}

\paragraph{Thermal expansion.}
Linear thermal expansion strain relative to $T_{\mathrm{ref}} = 293.15$ K:
\begin{equation}
  \varepsilon^{\mathrm{th}}(T) = \varepsilon_{\mathrm{mat}}(T) - \varepsilon_{\mathrm{mat}}(T_{\mathrm{ref}}).
\end{equation}

\paragraph{Hooke's law.}  In cylindrical coordinates, for hoop ($h$), radial
($r$), and axial ($z$) components:
\begin{align}
  E \, (\varepsilon_h - \varepsilon^{\mathrm{th}}) &= \sigma_h - \nu(\sigma_r + \sigma_z),\\
  E \, (\varepsilon_r - \varepsilon^{\mathrm{th}}) &= \sigma_r - \nu(\sigma_h + \sigma_z),\\
  E \, (\varepsilon_z - \varepsilon^{\mathrm{th}}) &= \sigma_z - \nu(\sigma_h + \sigma_r).
\end{align}
The same form applies to both fuel (subscript $f$) and cladding (subscript $c$)
with their respective moduli.

\paragraph{Strain compatibility.}  Between adjacent rings $i$ and $i+1$:
\begin{equation}
  \Delta r \, \bar{\varepsilon}_r + r_i \varepsilon_{h,i} - r_{i+1} \varepsilon_{h,i+1} = 0,
\end{equation}
where $\bar{\varepsilon}_r = (\varepsilon_{r,i} + \varepsilon_{r,i+1})/2$.

\paragraph{Radial stress equilibrium.}  Between adjacent rings:
\begin{equation}
  \Delta r \, \bar{\sigma}_h + r_i \sigma_{r,i} - r_{i+1} \sigma_{r,i+1} = 0.
\end{equation}

\paragraph{Boundary conditions (fuel).}
\begin{itemize}
  \item \emph{Inner surface (solid pellet):} symmetry, $\sigma_{r,0} = \sigma_{h,0}$.
  \item \emph{Outer surface, open gap:} $\sigma_{r,n_f-1} = -p_{\mathrm{gas}}$.
  \item \emph{Outer surface, closed gap:} axial strain continuity, $\varepsilon^f_z = \varepsilon^c_z$.
  \item \emph{Axial force balance:} $\sum_i A_i \sigma_{z,i} + $ gas/coolant pressure terms = 0.
\end{itemize}

\paragraph{Boundary conditions (cladding).}
\begin{itemize}
  \item \emph{Inner surface, open gap:} $\sigma_{r,0} = -p_{\mathrm{gas}}$.
  \item \emph{Inner surface, closed gap:} continuity of radial stress with fuel.
  \item \emph{Outer surface:} $\sigma_{r,n_c-1} = -p_{\mathrm{coolant}}$.
\end{itemize}

\paragraph{Deformed geometry.}  After solving for strains, the current node
radii and ring thicknesses are:
\begin{align}
  r_i &= r_{0,i}(1 + \varepsilon_{h,i}),\\
  \Delta r_i &= \Delta r_0 (1 + \bar{\varepsilon}_{r,i}),\\
  \Delta z &= \Delta z_0 (1 + \varepsilon_z).
\end{align}
The fuel outer radius $r_{\mathrm{fo}}$ and cladding inner radius
$r_{\mathrm{ci}}$ follow from the outermost node radii, giving the
instantaneous gap width:
\begin{equation}
  \delta_{\mathrm{gap}} = r_{\mathrm{ci}} - r_{\mathrm{fo}}.
\end{equation}

\subsection{DAE System and SUNDIALS IDA}

The complete set of equations per axial layer $j$ is a
differential-algebraic system (DAE):
\begin{equation}
  F(t, \mathbf{y}, \dot{\mathbf{y}}) = \mathbf{0},
\end{equation}
where $\mathbf{y}$ collects all state variables.  The total number of
equations per layer is
\begin{equation}
  n_{\mathrm{eq},j} = 6 + 7 n_f + 7 n_c.
\end{equation}
For $n_f=5$, $n_c=3$: $n_{\mathrm{eq},j} = 62$.

The system is solved by SUNDIALS IDA (Implicit Differential-Algebraic solver)
using:
\begin{itemize}
  \item BDF (backward differentiation formulae) time integration,
  \item Newton iteration for the nonlinear system at each time step,
  \item Dense direct linear solve (LAPACK via \texttt{SUNLinSol\_Dense}).
\end{itemize}

%===========================================================================
\section{Material Property Correlations}
%===========================================================================

\subsection{UO2 Fuel}

\paragraph{Thermal conductivity} (Philipponneau, J.~Nucl.~Mat.~188, 1992,
at zero Pu content):
\begin{equation}
  k(T, \rho, B) = \left[\frac{1}{a_c + 2.493\times10^{-4} T}
                        + 88.4\times10^{-12} T^3\right]
                  \frac{1-p}{1+2p},
\end{equation}
where $p = 1 - \rho/\rho_{\mathrm{th}}$ is the porosity,
$\rho_{\mathrm{th}} = 10960$~kg/m$^3$,
and $a_c = 1.320\sqrt{0.0093} - 0.091 + 3.8\times10^{-3} B/9.33$.

\paragraph{Heat capacity} (Popov et al.\ ORNL/TM-2000/351):
Tabular data from 300–3100~K, linearly interpolated.

\paragraph{Thermal expansion} (MATPRO):
\begin{equation}
  \varepsilon^{\mathrm{th}}(T) = 10^{-5} T - 3\times10^{-3} + 4\times10^{-2} e^{-5000/T}
                                 - \varepsilon^{\mathrm{th}}(293.15).
\end{equation}

\paragraph{Young's modulus} (MATPRO):
\begin{equation}
  E(T,\rho) = 2.334\times10^{11}\,(1 - 2.752\,p)\,(1 - 1.0915\times10^{-4}\,T) \cdot f_{\mathrm{soft}},
\end{equation}
with softening factor $f_{\mathrm{soft}}=\max(0, 1-[(T-T_0)/(T_m-T_0)]^{12})$
for $T > T_0 = T_m - 150$~K ($T_m = 3120$~K for UO$_2$).

\paragraph{Poisson's ratio} (MATPRO): $\nu = 0.276$.

\subsection{AIM1 Austenitic Stainless Steel}

Source: Luzzi et al., RdS/PAR2013/022.

\begin{align}
  k(T) &= 13.95 + 0.01163\,(T - 273.15) \quad [\mathrm{W/(m\cdot K)}],\\
  c_p(T) &= 431.0 + 0.177\,T + 8.72\times10^{-5}\,T^2 \quad [\mathrm{J/(kg\cdot K)}],\\
  E(T) &= 2.027\times10^{11} - 8.167\times10^{7}\,(T-273.15) \quad [\mathrm{Pa}],\\
  \nu &= 0.289.
\end{align}

Thermal expansion strain (integrated from the AIM1 coefficient):
\begin{equation}
  \varepsilon^{\mathrm{th}}(T) = (-0.2177\times10^{-2} + 6.735\times10^{-6}\,T
      + 5.12\times10^{-9}\,T^2 - 2.248\times10^{-12}\,T^3
      + 3.933\times10^{-16}\,T^4) - \varepsilon^{\mathrm{th}}(293.15).
\end{equation}

Meyer hardness (SAS4A correlation):
\begin{equation}
  H_M(T) = \begin{cases}
    5.961\times10^9 \, T^{-0.206} & T \le 893.9~\mathrm{K},\\
    2.750\times10^{28} \, T^{-6.530} & T > 893.9~\mathrm{K}.
  \end{cases}
\end{equation}

\subsection{Liquid Sodium Coolant}

Source: Fink \& Leibowitz, ANL/RE-95/2 (1995), valid 371–1156~K.

\begin{align}
  k(T) &= 124.67 - 0.11381\,T + 5.5226\times10^{-5}\,T^2 - 1.1842\times10^{-8}\,T^3,\\
  c_p(T) &= 1658.2 - 0.84305\,T + 4.5090\times10^{-4}\,T^2 - 2992658/T^2,\\
  \rho(T) &= 1011.8 - 0.22054\,T - 1.9226\times10^{-5}\,T^2 + 5.6371\times10^{-9}\,T^3.
\end{align}

%===========================================================================
\section{Python Interface}
%===========================================================================

The Python module \texttt{fred} is built as a shared library
(\texttt{fred.cpython-*.so}) and exposes the C\texttt{++} objects directly
in memory via \texttt{pybind11}.  No data copying between Python and C\texttt{++}
is required during time integration.

\begin{lstlisting}[language=Python]
import fred
import numpy as np

# Geometry: 5 fuel nodes, 3 clad nodes, 10 axial layers
geom = fred.FuelRodGeometry()
geom.nf, geom.nc, geom.nz = 5, 3, 10
geom.rfi0 = 0.0;  geom.rfo0 = 4.1e-3
geom.rci0 = 4.2e-3; geom.rco0 = 4.75e-3
geom.dz0  = [0.01] * 10
geom.vgp  = 1e-5
geom.build()

# Materials
fuel = fred.UO2(reference_density=10400.0)
clad = fred.AIM1(reference_density=7900.0)

# Solver
solver = fred.FredRodSolver(geom, fuel, clad)
solver.set_power_density_history([0, 100, 1e6], [0, 2e8, 2e8])
solver.set_coolant_temperature([0, 1e6], [700, 700])
solver.set_coolant_pressure(0.5)
solver.set_gas_inventory(1e-5)
solver.set_initial_temperature(700.0)
solver.set_tolerances(1e-6, 1e-8)

solver.run(tend=1e6, dtout=1e4)

times = solver.time_points()
T_peak = solver.peak_fuel_temperature()
\end{lstlisting}

Python-defined custom materials are also supported: subclass
\texttt{fred.FuelPelletMaterial} and override the virtual methods.

%===========================================================================
\section{Build Instructions}
%===========================================================================

\subsection{Dependencies}

\begin{itemize}
  \item GCC $\ge$ 11 or Clang $\ge$ 14 (C\texttt{++}17 required)
  \item SUNDIALS 7.x (IDA, dense linear solver) — already installed at
        \texttt{/home/chan\_y/Documents/fred/SUNDIALS/INSTDIR}
  \item Python 3.10+ with pybind11 (\texttt{pip install pybind11})
\end{itemize}

\subsection{Compilation}

\begin{lstlisting}[language=bash]
cd fred_platform
make                # builds fred_rod.x (C++ executable)
make pymodule       # builds fred.cpython-*.so (Python module)
\end{lstlisting}

The Python module is placed in \texttt{build/} and can be imported directly
with \texttt{import sys; sys.path.insert(0, 'build')}.

\end{document}
